\scanpage{31}
\mspage{16}

\uline{Proof} : (I) \quad By (6.3), there are indeed no more than two automorphisms.

(II) and (III) \quad When the curves are put in the form (2.1), it follows from (2.5) that the automorphisms are described by their action on $\pi$ and that $\lambda\in G_m$ defines an automorphism if and only if
\[
c_4=\lambda^4\,c_4 \qquad\text{and}\qquad c_6=\lambda^6\,c_6 \, ,
\]
whence the assertions.

(IV) \quad If the curve is put in the form (6.1), the automorphisms are given by (1.2) : $x'\handwritten{\mathrel{\tikz[baseline=-.32ex]{\draw[->,thin] (0,0) .. controls (.19em,.39ex) and (.28em,.6ex) .. (.57em,.18ex) .. controls (.78em,-.23ex) and (1.02em,-.08ex) .. (1.28em,0) -- (2.28em,.04ex);}}}x$, for $s=t=0$ and
\[
\handwritten{\left\{}\begin{aligned}
&u^4=1\\
&r^3+ra_4+a_6(1-u^2)=0
\end{aligned}\right.
\]
and the composition law by
\[
(u',r')\quad (u,r)=(u\,u',u^{\raisebox{-.38ex}{$\scriptstyle\prime$}\,2}r+r')
\]

When the curve is put in the form (6.7), these equations reduce to $u\in\mu_4$ and $r\in F_3$ .

(V) \quad If the curve is put in the form (6.3), the automorphisms are given by (1.2) : $x'\handwritten{\mathrel{\tikz[baseline=-.32ex]{\draw[->,thin] (0,0) .. controls (.14em,.37ex) and (.25em,.52ex) .. (.44em,.13ex) .. controls (.62em,-.25ex) and (.86em,-.2ex) .. (1.09em,.08ex) .. controls (1.3em,.31ex) and (1.52em,.23ex) .. (1.67em,.20ex) -- (2.36em,.16ex);}}}x$, for $r=s^2$ and
\[
\handwritten{\left\{}\begin{aligned}
&u^3=1\\
&s^4+sa_3+a_4(1-u)=0\\
&t^2+ta_3+s^6+s^2\,a_4=0
\end{aligned}\right.
\]%
% archive sheet number: 30
